% Part: proof-theory
% Chapter: natural-deduction
% Section: sequents

\documentclass[../../../include/open-logic-section]{subfiles}

\begin{document}

\olfileid{pt}{nat}{seq}

\olsection{استنتاج طبیعی با دنباله‌ها: \Log{N2c} و~\Log{N2i}}

!!^a{proof}~$\delta$ در \Log{N1c} یا \Log{N1i} نشان می‌دهد که
!!{formula} پایانی آن، یعنی~$!A$، از فرض‌های بازش به دست می‌آید. اگر
$\Gamma$ مجموعهٔ !!{formula}های~$!B$ باشد که $!B^x$ به‌عنوان فرض باز
در~$\delta$ ظاهر می‌شود، می‌توانیم این را به صورت
$\delta\Proves\Gamma\Sequent !A$ بنویسیم. این امر پیشنهاد می‌کند که
دستگاه استنتاج طبیعی را برای دنباله‌ها نیز صورت‌بندی کنیم. در چنین
دستگاهی، هر دنباله اطلاعات مربوط به فرض‌های بازی را حمل می‌کند که
!!{formula} سمت راست به آن‌ها وابسته است؛ یعنی فرض‌هایی که در
زیر-!!{proof} منتهی به آن دنباله بازند. این دستگاه، استنتاج طبیعی با
دنباله‌ها یا استنتاج طبیعی «به سبک دنباله‌ای» است و دستگاه‌های حاصل را
\Log{N2c} و~\Log{N2i} می‌نامیم.

برای نزدیک‌ترکردن تناظر میان \Log{N1} و \Log{N2}، مجموعه‌های~$\Gamma$
در سمت چپ را متشکل از !!{formula}های برچسب‌دار~$x:!B$ می‌گیریم. چون
قواعد \Log{N1} فقط بر !!{formula}های مقدمات، یعنی !!{formula}های پایانی
زیر-!!{proof}های بی‌واسطه، عمل می‌کنند، قواعد \Log{N2} نیز فقط بر تالی
دنباله عمل می‌کنند. بنابراین می‌توانیم فرمول‌های برچسب‌دار سمت چپ
دنباله‌ها را صرفاً \emph{بافت} بنامیم.

به جای فرض‌ها، بالاترین دنباله‌های برهان‌های \Log{N2} اصولی به شکل
\[x:!A \Sequent !A.\]
هستند. قواعد دقیقاً متناظر با قواعد \Log{N1}اند و در \olref{tab:N2}
آمده‌اند.

\subfile{rules-N2}

چون باید، همانند \Log{N1c} و \Log{N1i}، اجازه دهیم قواعد
\Intro{\lif}، \Intro{\lnot}، \Elim{\lor} و \Elim{\lexists} بتوانند
فرضی را مرخص کنند اما ملزم به آن نباشند، کاربردهای قواعد متناظر در
\Log{N2c} و~\Log{N2i} را حتی هنگامی درست می‌دانیم که !!{formula}های
بافتی متناظر~$!A$، $!B$ و~$!A(a)$ در بافت مقدمهٔ مربوط حضور نداشته
باشند.

\begin{prop}
اگر $\delta$ !!a{proof}ی در~\Log{N1c} یا \Log{N1i} برای~$!A$ باشد،
آنگاه در~\Log{N2c} یا به‌ترتیب \Log{N2i}، !!a{proof}~$\delta'$ای برای
$\Gamma\Sequent !A$ وجود دارد که در آن $\Gamma$ مجموعهٔ همهٔ
~$x:!B$هایی است که $!B^x$ فرض باز~$\delta$ است.
\end{prop}

\begin{proof}
  با استقرا بر $n=\pheight{\delta}$. اگر $n=0$، $\delta$ تنها فرض
  باز~$!A^x$ است. در این صورت $\Gamma=\{x:!A\}$ و
  $\Gamma\Sequent !A$ اصلی از \Log{N2c} یا~\Log{N2i} است.

  اگر $n>0$، حالت‌ها را برحسب آخرین استنتاج~$\delta$ جدا می‌کنیم. فقط
  یک حالت را بررسی می‌کنیم و بقیه را به‌عنوان تمرین می‌گذاریم.

  فرض کنید آخرین استنتاج \Intro{\lif} باشد. آنگاه $\delta$ چنین پایان
  می‌یابد:
  \[
  \AxiomC{$\Discharge{!B}{x}$}
  \RightLabel{$\delta_1$}
  \DeduceC{$!C$}
  \DischargeRule{\Intro{\lif}}{x}
  \UnaryInfC{$!B \lif !C$}
  \DisplayProof
  \]
  بنا بر فرض استقرا، در~\Log{N2c} یا~\Log{N2i} !!a{proof}~$\delta_1'$
  برای $\Gamma'\Sequent !C$ وجود دارد که $\Gamma'$ فرض‌های برچسب‌دار
  باز در~$\delta_1$ هستند. اگر $!B^x$ فرض باز~$\delta_1$ باشد، در
  استنتاج \Intro{\lif} مرخص می‌شود و فرض‌های باز~$\delta$ برابرند با
  $\Gamma=\Gamma'\setminus\{x:!B\}$. به بیان دیگر، $\delta_1'$
  !!a{proof}ی برای~$x:!B,\Gamma\Sequent !C$ است و می‌توانیم $\delta'$
  را چنین بگیریم:
  \[
  \AxiomC{}
  \RightLabel{$\delta_1'$}
  \Deduce$x:!B, \Gamma \fCenter !C$
  \RightLabel{\Intro{\lif}}
  \UnaryInf$\Gamma \fCenter !B \lif !C$
  \DisplayProof
  \]
  اگر $!B^x$ فرض باز~$\delta_1$ نباشد، \Intro{\lif} هیچ فرضی را مرخص
  نمی‌کند و فرض‌های باز $\delta$ و~$\delta_1$ یکسان‌اند. چون لازم نیست
  فرمول بافتی~$x:!B$ در مقدمهٔ \Intro{\lif} در~\Log{N2c} و
  ~\Log{N2i} حاضر باشد، می‌توانیم $\delta'$ را چنین بگیریم:
  \[
  \AxiomC{}
  \RightLabel{$\delta_1'$}
  \Deduce$\Gamma \fCenter !C$
  \RightLabel{\Intro{\lif}}
  \UnaryInf$\Gamma \fCenter !B \lif !C$
  \DisplayProof
  \]

  حالت‌هایی که $\delta$ به قاعده‌ای دیگر پایان می‌یابد مشابه‌اند.
\end{proof}

\begin{prop}
اگر $\delta$ !!a{proof}ی در \Log{N2c} یا \Log{N2i} برای
$\Gamma\Sequent !A$ باشد، آنگاه در \Log{N1c} یا به‌ترتیب \Log{N1i}،
!!a{proof}~$\delta'$ای با !!{formula} پایانی~$!A$ و فرض باز~$!B^x$
برای هر $x:!B\in\Gamma$ وجود دارد.
\end{prop}

\begin{proof}
  باز با استقرا بر !!{height}~$n$ برهان~$\delta$ پیش می‌رویم. اگر
  $n=0$، $\delta$ اصل $x:!A\Sequent !A$ است و !!{proof}~$\delta'$ همان
  فرض~$!A^x$ است.

  اگر $n>0$، حالت‌ها را برحسب آخرین استنتاج~$\delta$ جدا می‌کنیم. برای
  نمونه، اگر آن استنتاج \Intro{\lif} باشد، $\delta$ به یکی از دو شکل
  زیر پایان می‌یابد:
  \[
  \bottomAlignProof
  \AxiomC{}
  \RightLabel{$\delta_1$}
  \Deduce$x:!B, \Gamma \fCenter !C$
  \RightLabel{\Intro{\lif}}
  \UnaryInf$\Gamma \fCenter !B \lif !C$
  \DisplayProof
  \quad\text{ یا }\quad
  \bottomAlignProof
  \AxiomC{}
  \RightLabel{$\delta_1$}
  \Deduce$\Gamma \fCenter !C$
  \RightLabel{\Intro{\lif}}
  \UnaryInf$\Gamma \fCenter !B \lif !C$
  \DisplayProof
  \]
  بنا بر فرض استقرا، در~\Log{N1c} یا~\Log{N1i} !!a{proof}
  ~‌$\delta_1'$ای برای~$!C$ وجود دارد. در حالت نخست $!B^x$ فرض باز
  ~‌$\delta_1'$ است و در حالت دوم نیست. می‌توانیم !!{proof}~$\delta'$
  را به‌ترتیب چنین بگیریم:
  \[
  \bottomAlignProof
  \AxiomC{$\Discharge{!B}{x}$}
  \RightLabel{$\delta_1'$}
  \DeduceC{$!C$}
  \DischargeRule{\Intro{\lif}}{x}
  \UnaryInfC{$!B \lif !C$}
  \DisplayProof
  \quad\text{ یا }\quad
  \bottomAlignProof
  \AxiomC{}
  \RightLabel{$\delta_1'$}
  \DeduceC{$!C$}
  \RightLabel{\Intro{\lif}}
  \UnaryInfC{$!B \lif !C$}
  \DisplayProof
  \]
\end{proof}

\end{document}
